\documentclass[10pt,twocolumn]{article}
\usepackage[T1]{fontenc}
\usepackage[utf8]{inputenc}
\usepackage{lmodern}
\usepackage[margin=1.8cm,top=2.4cm,bottom=2cm,columnsep=0.6cm]{geometry}
\usepackage{fancyhdr}
\usepackage{titlesec}
\usepackage{booktabs}
\usepackage{amsmath,amssymb,amsfonts}
\usepackage{microtype}
\usepackage{xcolor}
\usepackage{mdframed}
\usepackage{parskip}
\usepackage{enumitem}
\usepackage{hyperref}

\definecolor{rulered}{RGB}{180,30,30}
\definecolor{boxgray}{RGB}{245,245,245}
\definecolor{keywordgray}{RGB}{100,100,100}

\pagestyle{fancy}
\fancyhf{}
\fancyhead[L]{\textsc{\small Claude Mag}}
\fancyhead[C]{\small\textit{Machine Learning Research Digest}}
\fancyhead[R]{\small 2026-09-30}
\fancyfoot[C]{\small\thepage}
\renewcommand{\headrulewidth}{0.8pt}

\titleformat{\section}{\normalsize\bfseries\color{rulered}}{}{0pt}{}%
  [\vspace{-4pt}\color{rulered}\rule{\columnwidth}{0.4pt}]
\titlespacing{\section}{0pt}{8pt}{4pt}

\hypersetup{colorlinks=true,urlcolor=rulered,linkcolor=black}

\begin{document}
\setlength{\parindent}{0pt}
\setlength{\parskip}{4pt}

{\small\color{keywordgray}\textbf{Keywords:}
  diffusion LLMs \textbullet{}
  KV caching \textbullet{}
  speculative decoding \textbullet{}
  inference acceleration}\par\vspace{4pt}

{\LARGE\bfseries\raggedright
  Flash-dLLM}\par\vspace{4pt}

{\large\itshape\raggedright
  IO-Aware KV Caching and Parallel Decoding for Fast,
  Memory-Efficient Diffusion LLMs}\par\vspace{6pt}

{\small\textbf{By} Quan Nguyen-Tri, Mukul Ranjan \&
  Zhiqiang Shen (VILA Lab, MBZUAI)}\par\vspace{2pt}

{\small\color{keywordgray}arXiv:2609.26796 \textbullet{} September 2026}
\par\vspace{2pt}\noindent{\color{rulered}\rule{\columnwidth}{0.6pt}}\vspace{6pt}

An IO-aware fused KV-cache kernel combined with a self-drafting decode strategy
eliminates the HBM bandwidth bottleneck that has crippled diffusion LLM
inference, delivering 5.1$\times$ and 11.0$\times$ speedups over the strongest
prior method on GSM8K and HumanEval without any model modification.

\begin{mdframed}[backgroundcolor=boxgray,linecolor=rulered,linewidth=1pt,
    innerleftmargin=6pt,innerrightmargin=6pt,
    innertopmargin=6pt,innerbottommargin=6pt]
\textbf{\small AT A GLANCE}\par\vspace{4pt}
\begin{description}[leftmargin=0pt,labelsep=4pt,itemsep=2pt,parsep=0pt,
    font=\small\bfseries]
  \item[Problem:] Diffusion LLMs require full KV-cache reads at every denoising
    step; GPU HBM I/O is the dominant bottleneck, not compute.
  \item[Why it matters:] dLLMs promise parallel token generation, but the IO
    bottleneck has prevented them from reaching autoregressive LLM speeds.
  \item[Method:] Fused IO-aware CUDA kernel reduces redundant HBM reads;
    the dLLM itself serves as drafter and verifier, eliminating auxiliary models.
  \item[Result:] 5.1$\times$ speedup on GSM8K and 11.0$\times$ on HumanEval
    vs.\ Elastic-Cache; improved memory efficiency at equal quality.
  \item[Evidence:] Mathematical reasoning and code generation benchmarks
    on A10G, A100, and L4 GPUs.
\end{description}
\end{mdframed}

\section{The Big Picture}

Diffusion language models (dLLMs) generate text by iteratively denoising
masked tokens rather than decoding one token at a time. Their non-autoregressive
structure should, in principle, allow far more parallelism than autoregressive
models. In practice, dLLM inference has remained slow because every denoising
step requires reading the entire KV cache---a cost proportional to sequence
length that must be paid at each of the $T$ iterations.

Flash-dLLM diagnosises this as an I/O bottleneck: on all tested GPUs, the
attention computation in a KV-cache-enabled dLLM is \emph{memory-bandwidth-bound},
not compute-bound. Arithmetic intensity (FLOPs per HBM byte) falls far below
the GPU's compute-to-memory ratio, meaning that the GPU arithmetic units are
idle most of the time, waiting for data.

\section{Why Existing Approaches Fall Short}

\textbf{Standard autoregressive KV caching} exploits left-to-right structure;
dLLMs access the full bidirectional context at every step, so per-step IO cost
is proportional to sequence length regardless of generation position.

\textbf{Elastic-Cache and token eviction} reduce cache size but preserve the
naive kernel structure: the remaining tokens still require an unfused
read-attend-write pipeline that is IO-inefficient.

\textbf{Auxiliary-model draft-and-verify} (as used in speculative decoding for
autoregressive models) adds a second model to load and run, increasing memory
and latency per step.

\textbf{Flash Attention} for autoregressive models was designed for compute-bound
settings; its tiling strategy does not generalise to the full-sequence,
bidirectional reads required by dLLMs.

\section{Core Mechanism}

\textbf{IO-Aware Fused Kernel.} Flash-dLLM fuses the KV-cache read, attention
computation, and output write into a single kernel. Intermediate values are
kept in on-chip SRAM, eliminating round-trips to HBM for the intermediate
attention maps. The tiling strategy is tuned for the bidirectional, full-sequence
access pattern of dLLMs rather than the causal sliding-window pattern of
autoregressive models.

\textbf{Self-Drafting Decode Strategy.} Early denoising steps operate at high
noise levels and produce coarse but fast predictions. Flash-dLLM uses these
as a draft: the KV cache computed in early steps is reused by later steps,
and positions where the early draft matches the final refined output avoid
recomputation. The dLLM itself is both drafter and verifier---no auxiliary
model is loaded.

\section{Technical Formulation}

Let $\mathbf{x}^{(t)}$ be the partially denoised sequence at step $t$,
$f_\theta$ the dLLM, $L$ the sequence length, $d$ the head dimension,
$H$ the number of heads, $l$ the number of layers.

Naive KV-cache IO cost per step:
\[
\mathrm{IO}(t) = 2 \cdot L \cdot d \cdot H \cdot l \cdot \mathrm{sizeof}(\mathrm{fp16})
\]

Fused-kernel IO cost (tile size $B_{\mathrm{SRAM}}$):
\[
\mathrm{IO}_{\mathrm{fused}}(t)
  = \frac{L \cdot d \cdot H \cdot l}{B_{\mathrm{SRAM}}} \cdot \mathrm{sizeof}(\mathrm{fp16})
\]

The draft-and-verify step amortises IO over $T_{\mathrm{draft}}$ steps: positions
where the draft $\hat{\mathbf{x}}^{(0)}$ at step $t$ matches the verified output
at step $t-1$ skip re-attention, contributing a further multiplicative reduction
in total IO.

\section{Learning or Inference Procedure}

Flash-dLLM requires no training and no model modification:
\begin{enumerate}[itemsep=2pt,parsep=0pt]
  \item Compile the IO-aware fused CUDA kernel for the target GPU architecture.
  \item Run the first $T_{\mathrm{draft}}$ denoising steps with the fused kernel;
    record draft token predictions from early steps.
  \item Verify drafts in the remaining $T - T_{\mathrm{draft}}$ steps; reuse
    KV cache for positions where draft is accepted.
  \item Return the fully denoised sequence.
\end{enumerate}

The kernel and draft parameters are tuned once per GPU model; no per-task
adaptation is needed.

\section{What the Guarantee Says}

The fused kernel is numerically exact: it computes the same attention values
as the unfused baseline up to floating-point rounding. Speedup is achieved
purely by reducing memory traffic, not by approximation. The self-drafting
strategy introduces an optimism bias (assuming later steps will agree with
earlier ones), but the verification step guarantees that accepted tokens are
correct.

\section{Experimental Findings}

\begin{itemize}[itemsep=2pt,parsep=0pt]
  \item Versus Elastic-Cache (strongest prior baseline):
    \textbf{5.1$\times$ speedup} on GSM8K (mathematical reasoning);
    \textbf{11.0$\times$ speedup} on HumanEval (code generation)
  \item Memory efficiency: lower peak HBM usage at equal or better quality
  \item Code benchmark shows larger gains because longer sequences and sparser
    denoising trajectories amplify IO savings
  \item Inference speed approaches autoregressive LLMs with Flash Attention
    at the same model size
\end{itemize}

\section{Ablations and Interpretation}

\textbf{Fused kernel alone} accounts for 3--5$\times$ of the speedup; the
self-drafting strategy contributes 1.5--2$\times$.

\textbf{Draft length $T_{\mathrm{draft}}$:} Shorter drafts reduce reuse and
speedup; longer drafts risk more corrections but still improve throughput
because IO savings outweigh correction costs.

\textbf{IO profile:} On all tested GPUs, attention accounts for over 70\% of
total inference time with KV caching, confirming the IO-bottleneck hypothesis.

\textbf{Arithmetic intensity:} Before the fused kernel, dLLM attention falls
below 10\% of the GPU roofline; after, it approaches 60--80\%, indicating
near-full utilisation of HBM bandwidth.

\section*{Reference}

Quan Nguyen-Tri, Mukul Ranjan, Zhiqiang Shen.
\textbf{Flash-dLLM: IO-Aware KV Caching and Parallel Decoding for Fast,
Memory-Efficient Diffusion LLMs}.
arXiv:2609.26796, September 2026.\\
\url{https://arxiv.org/abs/2609.26796}

\end{document}
